% 第 3 题程序框图：按固定源内容流矢量层精确值复刻（1 pt → 0.03515 cm，与全库同制）。
% 原点取框图竖轴 x=216.2232 pt、底缘（结束框底缘）y（顶向下）=509.4991 pt；y 向上（Y = 509.4991 − y_页）。
% 源实测（figure-inventory §1）：
%   7 框 —— 开始/结束＝圆角矩形 27.9988×16.8384 / 27.9988×17.0182（圆角半径 3.9819 横向 × 3.9858 纵向，
%   源为 4 段三次贝塞尔、控制点偏移 2.2000，本图取 rounded corners=0.1400cm 同制）；
%   k=1 矩形 32.0472×14.2905、k=k+1 矩形 50.6793×15.1659、k>4? 菱形 76.0365×18.8821（半对角线 38.0183×9.4410）、
%   S=sin(kπ/6) 矩形 57.2404×28.5967（内嵌真分式故高为数据框的 2 倍强；分数线在源中是独立描边
%   x 227.9894–240.1463 @ y=433.1499，本图由 \dfrac 自带分数线渲染）、
%   输出 S 平行四边形 50.8121×16.8579（斜边水平偏移 6.1351，上边偏右）。
%   连线 6 条竖直（竖轴 x=0，箭尖照源实测；源三角 2.9855 宽 × 3.9820 高，尖端向下）；
%   「否」支 3 段回折 —— 自菱形右顶点右行 (38.2997→52.4730)、竖直上行 (115.7156→162.2645)、
%   水平左行至竖轴 x=0，箭尖 (0,162.2645) 恰落在 k=1 框底与 k=k+1 框顶的正中点（源实测落点
%   y_页=347.2346 与两框间隙中点 347.2327 差 0.0019）——即回流到 k=k+1 的入线正中，
%   判定为当型循环；实心箭头共 7 枚（6 下 + 1 左）。全图无虚线、无圆形节点。
%   以下坐标均为本图坐标系（X = x_页 − 216.2232，Y = 509.4991 − y_页），与代码逐值一致。
%   框内文字与分支标签一律按源字形盒并集中心定位（figure-inventory §1.4）。
\begin{tikzpicture}[x=0.03515cm,y=0.03515cm,>=stealth,every node/.style={font=\small,inner sep=0pt}]
% ===== 7 框（外轮廓照源实测值）=====
\draw[line width=0.398bp,rounded corners=0.1400cm] (-13.9974,198.4148) rectangle (14.0014,215.2532); % 1 开始（圆角）
\draw[line width=0.398bp] (-16.0216,169.5524) rectangle (16.0256,183.8429);                            % 2 k=1
\draw[line width=0.398bp] (-25.3377,139.8146) rectangle (25.3416,154.9805);                            % 3 k=k+1
\draw[line width=0.398bp] (0,125.1567)--(38.0183,115.7156)--(0,106.2746)--(-38.0182,115.7156)--cycle;  % 4 菱形 k>4?
\draw[line width=0.398bp] (-28.6202,63.0200) rectangle (28.6202,91.6167);                              % 5 S=sin(kπ/6)
\draw[line width=0.398bp] (-25.4080,31.5901)--(-19.2729,48.4480)--(25.4041,48.4480)--(19.2690,31.5901)--cycle; % 6 输出 S
\draw[line width=0.398bp,rounded corners=0.1400cm] (-13.9974,0) rectangle (14.0014,17.0182);            % 7 结束（圆角）
% ===== 连线（6 竖直 + 1 回折；箭尖照源实测）=====
\draw[->,line width=0.398bp] (0,198.2155)--(0,184.0422);   % 开始→k=1
\draw[->,line width=0.398bp] (0,169.3531)--(0,155.1798);   % k=1→k=k+1（回折靶线）
\draw[->,line width=0.398bp] (0,139.6153)--(0,125.4420);   % k=k+1→菱形顶
\draw[->,line width=0.398bp] (0,105.9932)--(0,91.8199);    % 菱形底→S=sin（是支）
\draw[->,line width=0.398bp] (0,62.8245)--(0,48.6512);     % S=sin→输出 S
\draw[->,line width=0.398bp] (0,31.3908)--(0,17.2175);     % 输出 S→结束
\draw[->,line width=0.398bp] (38.2997,115.7156)--(52.4730,115.7156)--(52.4730,162.2645)--(0,162.2645); % 否支回折（回流 k=k+1 入线正中）
% ===== 框内文字与分支标签（照源字形盒并集中心定位）=====
\node at (-0.0002,207.4161) {开始};
\node at (-0.0047,175.6761) {\(k=1\)};
\node at (-0.0017,146.8131) {\(k=k+1\)};
\node at (-0.0007,114.8451) {\(k>4?\)};
\node at (-0.5981,76.3001) {\(S=\sin\dfrac{k\pi}{6}\)};
\node at (-0.3052,40.3496) {输出 \(S\)};
\node at (-0.0002,9.1181) {结束};
\node at (-8.9492,99.5931) {是};
\node at (38.2988,124.6861) {否};
\end{tikzpicture}
